\documentclass[10pt]{article}
\usepackage{graphicx}
\usepackage{tabularx}
\usepackage{fancyhdr}
\usepackage{hyperref}
\usepackage[
headheight=85pt,
headsep=10pt,
top=1.5in,
bottom=0.55in,
left=0.65in,
right=0.65in
]{geometry}
\usepackage{array}
\usepackage{enumitem}
\usepackage{booktabs}
\usepackage[T1]{fontenc}
\usepackage{mathptmx}
\usepackage{amsmath,amssymb}
\usepackage{xcolor}
\usepackage{subcaption}

\renewcommand{\headrulewidth}{1pt}
\setlength{\headheight}{85pt}
\setlength{\headsep}{10pt}

\fancyhead[L]{%
\begin{tabularx}{\linewidth}{@{}p{1.8cm}Xp{1.8cm}@{}}
\parbox[c][1.8cm][c]{1.8cm}{%
  \centering
  \includegraphics[width=1.8cm,height=1.6cm,keepaspectratio]{iit-madras-logo.png}%
}%
&
\parbox[c][1.8cm][c]{\linewidth}{%
  \centering
  {\Large\bfseries Touchlab: Weekly Research Report}\\[0.08cm]
  Applied Mechanics \& Biomedical Engineering, IIT Madras\\
  apm.iitm.ac.in/biomedical/touchlab\\[0.12cm]
  \textbf{From:} Sanath B S \quad\textbar\quad \textbf{Period:} $12^{\mathrm{th}}$ Nov 2026 -- $18^{\mathrm{th}}$ Nov 2026%
}%
&
\parbox[c][1.8cm][c]{1.8cm}{%
  \raggedleft
  \includegraphics[width=1.8cm,height=1.6cm,keepaspectratio]{tchlabLogo.PNG}%
}%
\end{tabularx}%
}

\pagestyle{fancy}

\begin{document}

% ═════════════════════════════════════════════════════════════════════════════
% PAGE 1: RESEARCHING REPOSITORIES, WORKFLOWS & ZERO-KEY IMAGE PIPELINES
% ═════════════════════════════════════════════════════════════════════════════
\section*{Researching Presentation Repositories, Workflows \& Visual Architectures}

\subsection*{1. Researching Presentation Frameworks \& Decoupled Agentic Pipelines}

This week, we researched open-source presentation engines, headless export runners, and layout compilation frameworks, including \textbf{Presenton}, \textbf{python-pptx}, and headless Chromium runtimes. Monolithic presentation tools that generate slides in a single prompt produce ungrounded, visually flat decks without teacher control. 

To solve this, we decoupled slide generation into a sequential 5-stage architecture (Figure~\ref{fig:architecture}):
\begin{itemize}[leftmargin=1.5em, itemsep=0.15em]
    \item \textbf{Input \& Scope:} Defines target topic, grade level, and curriculum passages.
    \item \textbf{Agentic Outline:} LLM extracts an editable sequence of concepts and formative checkpoints before rendering.
    \item \textbf{Layout \& Template Engine:} Maps outlines into validated Pydantic schema unions and design tokens.
    \item \textbf{Zero-Key Image Bridge:} Dynamically retrieves authentic schematics without external paid API keys.
    \item \textbf{Slide Studio \& Export:} Renders an interactive 16:9 canvas with AI co-pilot prompts and native PPTX/PDF export.
\end{itemize}

\vspace{0.1cm}

\begin{figure}[htbp]
\centering
\includegraphics[width=\textwidth,height=3.4cm,keepaspectratio]{fig_architecture.png}
\caption{\textbf{Decoupled Presentation Architecture:} End-to-end dataflow separating pedagogical outline extraction, layout compilation, zero-key image retrieval, and interactive studio rendering.}
\label{fig:architecture}
\end{figure}

\vspace{-0.3cm}
\subsection*{2. Researching Visual Pipelines, Diagram Sourcing \& Rate Limits}

A primary bottleneck in AI presentation tools is visual quality. Most engines rely on stock photo APIs (Pexels, Pixabay) with paused registrations, or free image endpoints with 45-second IP cooldowns that cause parallel requests to fail, leaving blank grey placeholder boxes. We developed a \textbf{zero-key multi-tier visual resolution bridge} to solve this:
\begin{itemize}[leftmargin=1.5em, itemsep=0.15em]
    \item \textbf{Domain Whitelisting:} Queries are routed to \textbf{DuckDuckGo Scientific Index} and \textbf{Wikimedia Commons Educational Repository}, strictly blacklisting commercial stock photo sites.
    \item \textbf{Local Normalization:} Diagrams are standardized to $1024\times768$ JPEGs via PIL, cached locally by prompt hash, and embedded into the PowerPoint binary package to guarantee 100\% offline-playable slides.
    \item \textbf{Concurrency Pacing:} Slides are generated with batch-size throttling ($\text{batch\_size} = 2$) to operate reliably within Google Gemini's 15 RPM free-tier limit without triggering HTTP 429 errors.
\end{itemize}

\vspace{0.1cm}

\begin{figure}[htbp]
\centering
\includegraphics[width=\textwidth,height=4.6cm,keepaspectratio]{fig_presenton_editor.png}
\caption{\textbf{Active Slide Studio Interface:} Live Presenton-OMNIX editing workspace showing direct in-canvas text manipulation, slide filmstrip navigation, and pedagogical configuration panel.}
\label{fig:presenton_editor}
\end{figure}

\clearpage

% ═════════════════════════════════════════════════════════════════════════════
% PAGE 2: MULTI-THEME 2x2 SLIDE PREVIEW MATRICES & VISUAL EVALUATION
% ═════════════════════════════════════════════════════════════════════════════
\subsection*{3. Multi-Theme Generated Slide Deck Outputs \& Visual Evaluation}

We validated our pipeline by generating complete 7-slide decks for Grade 9 Biology (\textit{"The Fundamental Unit of Life: Cell Structure and Functions"}) across three distinct design themes. Our visual engine eliminated blank boxes and watermarks, embedding authentic scientific diagrams across every slide with full typography and structured cards.

\vspace{0.15cm}

\begin{figure}[htbp]
\centering
\includegraphics[width=\textwidth,height=5.2cm,keepaspectratio]{fig_matrix_general.png}
\caption{\textbf{Theme 1: General (Visual Hero Panels):} 2$\times$2 matrix of rendered 16:9 slides featuring large rounded visual hero cards, fluid mosaic plasma membrane diagrams, and chloroplast plastid schematics.}
\label{fig:matrix_general}
\end{figure}

\vspace{-0.25cm}

\begin{figure}[htbp]
\centering
\begin{subfigure}[b]{0.485\textwidth}
    \centering
    \includegraphics[width=\textwidth,height=5.0cm,keepaspectratio]{fig_matrix_standard.png}
    \caption{\textbf{Theme 2: Standard (Technical Layout):} Structured side-by-side cards with mitochondria and cytology microscopy.}
\end{subfigure}
\hfill
\begin{subfigure}[b]{0.485\textwidth}
    \centering
    \includegraphics[width=\textwidth,height=5.0cm,keepaspectratio]{fig_matrix_modern.png}
    \caption{\textbf{Theme 3: Modern (Card Grid System):} Clean multi-card layout with vector iconography and concept descriptions.}
\end{subfigure}
\caption{\textbf{Cross-Theme Comparison:} 2$\times$2 slide preview matrices comparing the Standard technical theme against the Modern card grid theme on the same biology curriculum units.}
\label{fig:multi_theme_matrices}
\end{figure}

\vspace{-0.25cm}
\subsection*{Summary \& Immediate Next Steps}

\begin{itemize}[leftmargin=1.5em, itemsep=0.2em]
    \item Pre-bundle a curated library of standardized NCERT and OpenStax textbook diagrams for guaranteed label accuracy.
    \item Implement an in-canvas 1-click diagram swapper modal allowing teachers to replace visuals directly in the studio.
    \item Expose native OMNIX FastAPI proxy endpoints for direct slide generation and PPTX/PDF downloads.
\end{itemize}

\end{document}
